\section{Check if an Array Represents a Heap}
\label{sec:heap-check-array}

\problemstatement{Given an array $\textit{arr}[0..n-1]$, determine whether
it represents a valid \textbf{max-heap} when interpreted as a complete
binary tree in level order (every parent at index $i$ should be $\ge$ its
children at $2i+1$ and $2i+2$).}

\begin{patternbox}
This is a direct test of the array-as-tree arithmetic from
Section~\ref{sec:heap-implement}: the keyword is simply \emph{``does this
array, read as a complete binary tree, satisfy the heap invariant
everywhere.''} There is no clever algorithm to discover -- the entire
exercise is checking the invariant at every internal node using the
$\textit{left}(i) = 2i+1$, $\textit{right}(i) = 2i+2$ formulas directly.
\end{patternbox}

\subsection*{Understanding the problem carefully}

A leaf node (one with no children, i.e.\ $2i+1 \ge n$) trivially satisfies
the heap property, since it has nothing to compare against. So we only need
to check internal nodes -- indices $i$ with $2i+1 < n$ -- and verify that
$\textit{arr}[i] \ge \textit{arr}[2i+1]$ (when the left child exists) and
$\textit{arr}[i] \ge \textit{arr}[2i+2]$ (when the right child exists). If
every internal node passes both checks, the whole array is a valid
max-heap, because the heap property is purely local: validity at every
parent-child pair, checked independently, is both necessary and sufficient
for the global invariant to hold (there is no broader ordering constraint
a heap is required to satisfy beyond these local pairwise comparisons).

\begin{ideabox}[Why Checking Only Internal Nodes Suffices]
The heap property is defined entirely in terms of parent-versus-children
comparisons. A node with no children imposes no constraint of its own, so
the full set of constraints to verify is exactly the set of (parent, child)
pairs that exist in the array -- checking every internal node against its
existing children covers all of them exactly once.
\end{ideabox}

\subsection*{Algorithm}

\begin{enumerate}
  \item For $i \gets 0$ to $n-1$, while $2i+1 < n$ (i.e.\ $i$ is internal):
  \begin{enumerate}
    \item If $2i+1 < n$ and $\textit{arr}[i] < \textit{arr}[2i+1]$: return
          \texttt{false}.
    \item If $2i+2 < n$ and $\textit{arr}[i] < \textit{arr}[2i+2]$: return
          \texttt{false}.
  \end{enumerate}
  \item If every internal node passes, return \texttt{true}.
\end{enumerate}

\begin{algorithm}[H]
\caption{Check Array Represents a Max-Heap}
\begin{algorithmic}[1]
\Procedure{IsHeap}{\textit{arr}}
  \For{$i \gets 0$ to $n-1$}
    \State $\textit{left} \gets 2i+1$, $\textit{right} \gets 2i+2$
    \If{$\textit{left} < n$ \textbf{and} $\textit{arr}[i] < \textit{arr}[\textit{left}]$}
      \State \Return \textbf{false}
    \EndIf
    \If{$\textit{right} < n$ \textbf{and} $\textit{arr}[i] < \textit{arr}[\textit{right}]$}
      \State \Return \textbf{false}
    \EndIf
  \EndFor
  \State \Return \textbf{true}
\EndProcedure
\end{algorithmic}
\end{algorithm}

\subsection*{Dry run}

\dryrun{Take $\textit{arr} = [90, 15, 10, 7, 12, 2, 7, 3]$.}

\begin{itemize}
  \item $i=0$ (value $90$): children at $1,2$ are $15,10$. $90 \ge 15$ and
        $90 \ge 10$. OK.
  \item $i=1$ (value $15$): children at $3,4$ are $7,12$. $15 \ge 7$ and
        $15 \ge 12$. OK.
  \item $i=2$ (value $10$): children at $5,6$ are $2,7$. $10 \ge 2$ and
        $10 \ge 7$. OK.
  \item $i=3$ (value $7$): child at $7$ is $3$ (no right child, index $8$
        out of range). $7 \ge 3$. OK.
  \item $i=4,5,6,7$: no children ($2i+1 \ge 8$). Trivially OK.
\end{itemize}

Every internal node passes, so the array is a valid max-heap.

\subsection*{C++ implementation}

\begin{lstlisting}
#include <bits/stdc++.h>
using namespace std;

bool isHeap(vector<int>& arr) {
    int n = arr.size();
    for (int i = 0; i < n; i++) {
        int left = 2 * i + 1, right = 2 * i + 2;
        if (left < n && arr[i] < arr[left]) return false;
        if (right < n && arr[i] < arr[right]) return false;
    }
    return true;
}

int main() {
    vector<int> arr = {90, 15, 10, 7, 12, 2, 7, 3};
    cout << (isHeap(arr) ? "true" : "false") << endl;
    return 0;
}
\end{lstlisting}

\begin{complexitybox}
\textbf{Time Complexity.} Each index is visited once with constant work,
giving
\[
  O(n).
\]
\textbf{Space Complexity.} $O(1)$ auxiliary space.
\end{complexitybox}

\begin{takeawaybox}
The heap property is purely local (parent versus its own children), never
global -- a heap is not required to be sorted in any stronger sense. This
is precisely why heap operations are $O(\log n)$ rather than $O(n \log n)$:
restoring a single local violation after an insertion or removal never
requires re-examining the whole structure.
\end{takeawaybox}
